% 条件独立与条件化后出现的依赖。教学数据：
% (a) P(Z=0)=P(Z=1)=1/2；给定 Z=z，X,Y 独立同分布于 Bernoulli(p_z)，p_0=0.2，p_1=0.8。
%     条件联合表为乘积；混合后 P(X=1,Y=1)=0.34，而 P(X=1)P(Y=1)=0.25。
% (b) X,Y 为独立公平比特，Z=X xor Y；无条件联合表各格 0.25，给定 Z=0 或 Z=1 后只剩两格各 0.5。
% 格子底色深浅与概率成正比（FigureBlue 与白色按概率百分比混合）。
\begin{tikzpicture}[line width=.6pt,font=\fontsize{8}{10}\selectfont,>=Stealth,
  cell/.style={draw=FigureStroke,minimum width=9mm,minimum height=6mm,inner sep=0pt},
  note/.style={FigureMuted,font=\fontsize{7.5}{9}\selectfont}]
\foreach \ox/\oy/\ttl/\aa/\fa/\ab/\fb/\ba/\fc/\bb/\fd in {
  0/0/{给定$Z=0$}/0.64/64/0.16/16/0.16/16/0.04/4,
  36/0/{给定$Z=1$}/0.04/4/0.16/16/0.16/16/0.64/64,
  72/0/{按$Z$平均}/0.34/34/0.16/16/0.16/16/0.34/34,
  0/-35/{无条件}/0.25/25/0.25/25/0.25/25/0.25/25,
  36/-35/{给定$Z=0$}/0.5/50/0/0/0/0/0.5/50,
  72/-35/{给定$Z=1$}/0/0/0.5/50/0.5/50/0/0}{
  \begin{scope}[shift={(\ox mm,\oy mm)}]
  \node[FigureInk] at (9mm,11mm) {\ttl};
  \node[FigureMuted] at (4.5mm,5.5mm) {$y=0$};
  \node[FigureMuted] at (13.5mm,5.5mm) {$y=1$};
  \node[FigureMuted,anchor=east] at (-1mm,0) {$x=0$};
  \node[FigureMuted,anchor=east] at (-1mm,-6mm) {$x=1$};
  \node[cell,fill=FigureBlue!\fa!white] at (4.5mm,0) {$\aa$};
  \node[cell,fill=FigureBlue!\fb!white] at (13.5mm,0) {$\ab$};
  \node[cell,fill=FigureBlue!\fc!white] at (4.5mm,-6mm) {$\ba$};
  \node[cell,fill=FigureBlue!\fd!white] at (13.5mm,-6mm) {$\bb$};
  \end{scope}}
\node[note,anchor=west] at (92mm,0) {边缘均为$0.5$};
\node[note,anchor=west] at (92mm,-6mm) {$0.34\neq0.5\times0.5$};
\node[FigureInk,anchor=west] at (-12mm,-14mm) {(a) 给定类别$Z$时两表都是乘积，平均后不再是乘积};
\node[note,anchor=west] at (92mm,-35mm) {$Z=X\oplus Y$};
\node[FigureInk,anchor=west] at (-12mm,-49mm) {(b) 独立的两个比特在给定$Z$后互相确定};
\end{tikzpicture}
